\documentclass[10pt,a4paper]{amsart}
\usepackage{amsmath,amssymb,amsthm,graphicx}
\usepackage{verbatim}
\newenvironment{reference}{\comment}{\endcomment}
%\newenvironment{reference}{}{}
\newenvironment{slogan}{\comment}{\endcomment}
\newenvironment{history}{\comment}{\endcomment}

% For commutative diagrams we use Xy-pic
\usepackage[all]{xy}

% We use 2cell for 2-commutative diagrams.
\xyoption{2cell}
% Theorem environments.
%
\theoremstyle{plain}
\newtheorem{theorem}[subsection]{Theorem}
\newtheorem{proposition}[subsection]{Proposition}
\newtheorem{lemma}[subsection]{Lemma}

\theoremstyle{definition}
\newtheorem{definition}[subsection]{Definition}
\newtheorem{example}[subsection]{Example}
\newtheorem{exercise}[subsection]{Exercise}
\newtheorem{situation}[subsection]{Situation}

\theoremstyle{remark}
\newtheorem{remark}[subsection]{Remark}
\newtheorem{remarks}[subsection]{Remarks}

\numberwithin{equation}{subsection}

% Macros
%
\def\lim{\mathop{\mathrm{lim}}\nolimits}
\def\colim{\mathop{\mathrm{colim}}\nolimits}
\def\Spec{\mathop{\mathrm{Spec}}}
\def\Hom{\mathop{\mathrm{Hom}}\nolimits}
\def\Ext{\mathop{\mathrm{Ext}}\nolimits}
\def\SheafHom{\mathop{\mathcal{H}\!\mathit{om}}\nolimits}
\def\SheafExt{\mathop{\mathcal{E}\!\mathit{xt}}\nolimits}
\def\Sch{\mathit{Sch}}
\def\Mor{\mathop{\mathrm{Mor}}\nolimits}
\def\Ob{\mathop{\mathrm{Ob}}\nolimits}
\def\Sh{\mathop{\mathit{Sh}}\nolimits}
\def\NL{\mathop{N\!L}\nolimits}
\def\CH{\mathop{\mathrm{CH}}\nolimits}
\def\proetale{{pro\text{-}\acute{e}tale}}
\def\etale{{\acute{e}tale}}
\def\QCoh{\mathit{QCoh}}
\def\Ker{\mathop{\mathrm{Ker}}}
\def\Im{\mathop{\mathrm{Im}}}
\def\Coker{\mathop{\mathrm{Coker}}}
\def\Coim{\mathop{\mathrm{Coim}}}

% Boxtimes
%
\DeclareMathSymbol{\boxtimes}{\mathbin}{AMSa}{"02}

%
% Macros for moduli stacks/spaces
%
\def\QCohstack{\mathcal{QC}\!\mathit{oh}}
\def\Cohstack{\mathcal{C}\!\mathit{oh}}
\def\Spacesstack{\mathcal{S}\!\mathit{paces}}
\def\Quotfunctor{\mathrm{Quot}}
\def\Hilbfunctor{\mathrm{Hilb}}
\def\Curvesstack{\mathcal{C}\!\mathit{urves}}
\def\Polarizedstack{\mathcal{P}\!\mathit{olarized}}
\def\Complexesstack{\mathcal{C}\!\mathit{omplexes}}
% \Pic is the operator that assigns to X its picard group, usage \Pic(X)
% \Picardstack_{X/B} denotes the Picard stack of X over B
% \Picardfunctor_{X/B} denotes the Picard functor of X over B
\def\Pic{\mathop{\mathrm{Pic}}\nolimits}
\def\Picardstack{\mathcal{P}\!\mathit{ic}}
\def\Picardfunctor{\mathrm{Pic}}
\def\Deformationcategory{\mathcal{D}\!\mathit{ef}}
\usepackage{cleveref}
\crefname{originallemma}{original lemma}{original lemmas}
\Crefname{originallemma}{Original lemma}{Original lemmas}
\makeatletter
\newlabel{decent-spaces-lemma-filter-quasi-compact-quasi-separated}{{07ST}{source}{}{sourcepointer-1}{}}
\newlabel{decent-spaces-lemma-filter-quasi-compact-quasi-separated@cref}{{[originallemma][0][]07ST}{[1][0][]source}{}{}{}}
\makeatother
\crefname{originalreference}{original reference}{original references}
\Crefname{originalreference}{Original reference}{Original references}
\makeatletter
\newlabel{lemma-induction-principle}{{08DR}{source}{}{sourcepointer-2}{}}
\newlabel{lemma-induction-principle@cref}{{[originalreference][0][]08DR}{[1][0][]source}{}{}{}}
\makeatother
\crefname{originalreference}{original reference}{original references}
\Crefname{originalreference}{Original reference}{Original references}
\makeatletter
\newlabel{lemma-induction-principle-separated}{{08GQ}{source}{}{sourcepointer-3}{}}
\newlabel{lemma-induction-principle-separated@cref}{{[originalreference][0][]08GQ}{[1][0][]source}{}{}{}}
\makeatother
\crefname{originalreference}{original reference}{original references}
\Crefname{originalreference}{Original reference}{Original references}
\makeatletter
\newlabel{lemma-make-more-elementary-distinguished-squares}{{08GN}{source}{}{sourcepointer-4}{}}
\newlabel{lemma-make-more-elementary-distinguished-squares@cref}{{[originalreference][0][]08GN}{[1][0][]source}{}{}{}}
\makeatother
\title{Induction principle starting from a nonempty base subspace}
\author{The Stacks Project Authors}
\date{stacks-project:spaces-perfect.tex:local-snapshot-20261012; GNU FDL 1.2 or later}
\begin{document}
\maketitle
\noindent\textit{Editorial note.} All delivered mathematical text is the author's own, copied from the pinned Stacks Project chapter file spaces-perfect.tex; no statement or proof step was rewritten, reordered or removed.\par\medskip
\noindent\textit{Editorial note.} The counted proof is the author's complete argument: the proof is complete, with the finite affine refinement of the square handled in full\par\medskip
\noindent\textit{Editorial note.} The label lemma-induction-principle-enlarge names the counted statement, and the proof environment immediately following it in the source is delivered verbatim as its proof.\par\medskip
\noindent\textit{Editorial note.} The counted statement is carried by the author's own lemma environment in the chapter source.\par\medskip
\noindent\textit{Allowed standard prerequisites (editorial).}\begin{itemize}
\item elementary distinguished squares of algebraic spaces
\item the filtration of a quasi-compact quasi-separated space by quasi-compact opens from the decent spaces chapter
\item refining a distinguished square by an affine open cover so the square hypothesis applies
\item descending induction on the filtration
\end{itemize}\medskip
\ifdefined\hypertarget\hypertarget{sourcepointer-1}{}\fi
\noindent\textit{Source reference.} Original lemma 07ST: sources/stacks-chow/tags.txt, source lines 9913--9913. Prerequisite pointer invoked at proof line(s) 1738; it names a standard fact of a sibling chapter that the counted argument uses, and every step of the counted proof built on it is retained verbatim in the delivered body.\par\medskip
\ifdefined\hypertarget\hypertarget{sourcepointer-2}{}\fi
\noindent\textit{Source reference.} Original reference 08DR: sources/stacks-spaces-perfect/pinned-original-source.tex, source lines 1563--1563. Prerequisite pointer invoked at proof line(s) 1731; it names a standard fact of the same chapter that the counted argument uses, and every step of the counted proof built on it is retained verbatim in the delivered body.\par\medskip
\ifdefined\hypertarget\hypertarget{sourcepointer-3}{}\fi
\noindent\textit{Source reference.} Original reference 08GQ: sources/stacks-spaces-perfect/pinned-original-source.tex, source lines 1643--1643. Prerequisite pointer invoked at proof line(s) 1729; it names a standard fact of the same chapter that the counted argument uses, and every step of the counted proof built on it is retained verbatim in the delivered body.\par\medskip
\ifdefined\hypertarget\hypertarget{sourcepointer-4}{}\fi
\noindent\textit{Source reference.} Original reference 08GN: sources/stacks-spaces-perfect/pinned-original-source.tex, source lines 1541--1541. Prerequisite pointer invoked at proof line(s) 1769; it names a standard fact of the same chapter that the counted argument uses, and every step of the counted proof built on it is retained verbatim in the delivered body.\par\medskip
\begin{lemma}
\label{lemma-induction-principle-enlarge}
Let $S$ be a scheme. Let $X$ be a quasi-compact and quasi-separated
algebraic space over $S$. Let $W \subset X$ be a quasi-compact open
subspace. Let $P$ be a property of quasi-compact open subspaces of $X$.
Assume that
\begin{enumerate}
\item $P$ holds for $W$, and
\item for every elementary distinguished square
$(W_1 \subset W_2, f : V \to W_2)$ where 
such that
\begin{enumerate}
\item $W_1$, $W_2$ are quasi-compact open subspaces of $X$,
\item $W \subset W_1$,
\item $V$ is affine, and
\item $P$ holds for $W_1$,
\end{enumerate}
then $P$ holds for $W_2$.
\end{enumerate}
Then $P$ holds for $X$.
\end{lemma}

\begin{proof}
We can deduce this from Lemma \ref{lemma-induction-principle-separated},
but instead we will give a direct argument by explicitly redoing the proof of
Lemma \ref{lemma-induction-principle}. We will use the filtration
$$
\emptyset = U_{n + 1} \subset
U_n \subset U_{n - 1} \subset \ldots \subset U_1 = X
$$
and the morphisms $f_p : V_p \to U_p$ of
Decent Spaces, Lemma
\ref{decent-spaces-lemma-filter-quasi-compact-quasi-separated}.
We will prove that $P$ holds for $W_p = W \cup U_p$ by descending
induction on $p$. This will finish the proof as $W_1 = X$.
Note that $P$ holds for $W_{n + 1} = W \cap U_{n + 1} = W$
by (1). Assume $P$ holds for $W_{p + 1}$. Observe that
$W_p \setminus W_{p + 1}$ (with reduced induced subspace structure)
is a closed subspace of $U_p \setminus U_{p + 1}$.
Since $(U_{p + 1} \subset U_p, f_p : V_p \to U_p)$ is an elementary
distinguished square, the same is true for
$(W_{p + 1} \subset W_p, f_p : V_p \to W_p)$.
However (2) may not apply as $V_p$ may not be affine.
However, as $V_p$ is a quasi-compact scheme we may choose
a finite affine open covering $V_p = V_{p, 1} \cup \ldots \cup V_{p, m}$.
Set $W_{p, 0} = W_{p + 1}$ and
$$
W_{p, i} = W_{p + 1} \cup f_p(V_{p, 1} \cup \ldots \cup V_{p, i})
$$
for $i = 1, \ldots, m$. These are quasi-compact open subspaces of $X$
containing $W$. Then we have
$$
W_{p + 1} = W_{p, 0} \subset
W_{p, 1} \subset \ldots \subset
W_{p, m} = W_p
$$
and the pairs
$$
(W_{p, 0} \subset W_{p, 1}, f_p|_{V_{p, 1}}),
(W_{p, 1} \subset W_{p, 2}, f_p|_{V_{p, 2}}),\ldots,
(W_{p, m - 1} \subset W_{p, m}, f_p|_{V_{p, m}})
$$
are elementary distinguished squares by
Lemma \ref{lemma-make-more-elementary-distinguished-squares}.
Now (2) applies to each of these and we inductively
conclude $P$ holds for $W_{p, 1}, \ldots, W_{p, m} = W_p$.
\end{proof}

\end{document}
